\section{Experiments}
\label{sec:experiments}

We evaluate our semantic correspondence pipeline on standard benchmarks.
Our experiments aim to assess (i) the quality of different foundation models for correspondence, (ii) the impact of fine-tuning, (iii) the impact of soft-argmax refinement, and (iv) generalization across datasets.

%-------------------------------------------------------------------------

\paragraph{Experimental Setup}

We compare three vision foundation models: DINOv2 ViT-B/14~\cite{oquab2023dinov2}, DINOv3 ViT-B/16~\cite{simeoni2025dinov3}, and SAM ViT-B~\cite{kirillov2023segment}.
All models are evaluated using dense feature matching based on cosine similarity with a global argmax, then refined via a window-based soft-argmax.
Fine-tuning is performed using a contrastive loss.

%-------------------------------------------------------------------------

\paragraph{Datasets}

We conduct experiments on the SPair-71k dataset~\cite{min2019spair}, which provides dense semantic keypoint annotations across diverse object categories.
As an extension project, we additionally evaluate generalization on the PF-Willow dataset, testing cross-domain robustness without retraining.

%-------------------------------------------------------------------------

\paragraph{Evaluation Metric}

We evaluate semantic correspondence performance using the Percentage of Correct Keypoints (PCK), a standard metric for semantic matching.
Given a ground-truth keypoint in the target image, a predicted correspondence is considered correct if its Euclidean distance from the ground-truth location is smaller than a fraction $\alpha$ of the image size.
Following common practice, distances are normalized by the maximum of the image height and width.
We report PCK at multiple thresholds, $\alpha \in \{0.05, 0.1, 0.15, 0.2\}$, to capture performance under varying localization tolerances.

PCK is first computed at the \emph{per-keypoint} level by averaging correctness over all annotated keypoints in the dataset.
To enable a more fine-grained analysis, we additionally report \emph{per-category} PCK by aggregating keypoint-level scores within each object class.
This analysis highlights differences in correspondence accuracy across object categories with varying degrees of deformation and visual complexity.

Finally, we report \emph{per-image} PCK statistics by computing, for each image pair, the fraction of correctly matched keypoints and then aggregating these values across the dataset.
For per-image evaluation, we report the mean, median, standard deviation, minimum, and maximum PCK, providing insight into both average performance and variability across image pairs.

%-------------------------------------------------------------------------

\paragraph{Implementation Details}

Input images are resized according to the backbone requirements: $518\times518$ for DINOv2, $512\times512$ for DINOv3 and for SAM.
DINO models use ImageNet normalization, while SAM employs its official normalization scheme.

Fine-tuning is conducted by unfreezing the last two transformer layers and the normalization layers for DINO models, while for SAM we unfreeze the last two layers and the neck component.
All models are fine-tuned for a small number of epochs using a contrastive loss with a learning rate of $5e-6$, a temperature of $0.07$ and a batch size of $16$.

Matching is performed using cosine similarity between dense features, followed by a global argmax.
For refinement, we apply a window-based soft-argmax with temperature $T=0.01$ and window radius $W=5$.



%-------------------------------------------------------------------------

\subsection{Evaluation of Models without Soft-Argmax}

We first evaluate the correspondence performance of pretrained visual foundation models in a training-free setting, using a standard argmax-based
prediction to establish a strong baseline. Specifically, we compare DINOv2, DINOv3, and SAM by extracting frozen dense features from the final backbone
layer and performing correspondence via cosine similarity.

Table~\ref{tab:baseline_pck} reports the overall per-keypoint PCK scores at multiple thresholds.
Both DINO-based models significantly outperform SAM across all evaluation thresholds. DINOv2 achieves the best baseline performance, reaching 54.16\% PCK@0.1
and 70.78\% PCK@0.2, while DINOv3 shows comparable but slightly lower accuracy. These results indicate that self-supervised ViT models trained with
contrastive objectives naturally encode strong semantic alignment cues, even without any task-specific adaptation.

In contrast, SAM exhibits substantially lower correspondence accuracy, with only 23.24\% PCK@0.1 and 37.76\% PCK@0.2. This gap is consistent across
thresholds and highlights the limitations of segmentation-oriented representations when directly applied to dense semantic correspondence without
architectural or training adaptations.

Per-image statistics, reported in Table~\ref{tab:baseline_stats}, further confirm these trends.
DINO-based models show higher mean and median PCK values, as well as lower variance compared to SAM, indicating more stable correspondence
predictions across image pairs.

Overall, these results establish DINOv2 and DINOv3 as strong baselines for semantic correspondence and motivate further improvements through fine-tuning and enhanced prediction strategies.


\begin{table}[t]
    \centering
    \caption{\textbf{Overall per-keypoint PCK (\%) on SPair-71k using pretrained models without soft-argmax.}}
    \label{tab:baseline_pck}
    \setlength{\tabcolsep}{2pt}
    \begin{tabular}{lcccc}
    \toprule
    \textbf{Model} & \textbf{PCK@0.05} & \textbf{PCK@0.1} & \textbf{PCK@0.15} & \textbf{PCK@0.2} \\
    \midrule
    DINOv2 ViT-B/14 & \textbf{36.75} & \textbf{54.16} & \textbf{64.07} & \textbf{70.78} \\
    DINOv3 ViT-B/16 & 35.86 & 54.05 & 63.63 & 69.55 \\
    SAM ViT-B       & 13.15 & 23.24 & 31.03 & 37.76 \\
    \bottomrule
    \end{tabular}
\end{table}


\begin{table}[t]
    \centering
    \caption{\textbf{Per-image PCK@0.1 statistics (\%) on SPair-71k without soft-argmax.}}
    \label{tab:baseline_stats}
    \begin{tabular}{lccccc}
    \toprule
    \textbf{Model} & \textbf{Mean} & \textbf{Median} & \textbf{Std} & \textbf{Min} & \textbf{Max} \\
    \midrule
    DINOv2 ViT-B/14 & \textbf{51.43} & 50.00 & 30.30 & 0.00 & 100.00 \\
    DINOv3 ViT-B/16 & 48.54 & 50.00 & 30.92 & 0.00 & 100.00 \\
    SAM ViT-B       & 19.97 & 15.38 & 21.38 & 0.00 & 100.00 \\
    \bottomrule
    \end{tabular}
\end{table}



%-------------------------------------------------------------------------

\subsection{Fine-tuned models to beat the baselines}


\paragraph{Training Dynamics and Optimization Analysis}

During contrastive fine-tuning tests, we observed the optimization behavior of DINOv2.
The training loss decreases smoothly without signs of instability, while the validation PCK@0.1 increases
monotonically throughout training from approximately 61\% reaching a peak of 72.43\%.
These trends indicate stable optimization and convergence, despite fine-tuning only a small subset of model parameters.
The absence of divergence between training and validation loss curves also suggests that the model doesn't overfit
and that fine-tuning effectively improves semantic correspondence performance.

%-------------------------------------------------------------------------

\subsubsection{Layer Selection Analysis}

We analyze the impact of feature extraction depth by evaluating intermediate transformer layers for both DINOv2 and DINOv3.
Table~\ref{tab:intermediate_layers} reports per-keypoint PCK results for selected layers.

For both models, performance consistently improves when moving toward deeper layers.
In particular, the final layer achieves the highest PCK scores for both DINOv2 and DINOv3; while intermediate layers (\eg layer 6) perform poorly,
indicating that semantic correspondence emerges predominantly in higher-level representations.

This confirms that high-level semantic representations are crucial for accurate correspondence estimation.


\begin{table}
    \centering
    \caption{
        \textbf{Layer selection analysis for DINOv2 and DINOv3 on SPair-71k.}
        \textit{Top: DINOv2. Bottom: DINOv3}.
        Per-keypoint PCK results for selected intermediate transformer layers.
        Performance improves with depth, with the final layer achieving the highest scores,
        highlighting the importance of high-level semantic features for correspondence.}
    \label{tab:intermediate_layers}
    \begin{tabular}{lccc}
    \toprule
    \textbf{Layer} & \textbf{PCK@0.05} & \textbf{PCK@0.1} & \textbf{PCK@0.2} \\
    \midrule
    Layer 6  & 3.72\%  & 11.27\% & 25.25\% \\
    Layer 9  & 8.10\%  & 22.38\% & 40.16\% \\
    Layer 10 & 11.40\% & 31.80\% & 54.98\% \\
    Final    & \textbf{12.96\%} & \textbf{35.96\%} & \textbf{62.09\%} \\
    \bottomrule
    \toprule
    \textbf{Layer} & \textbf{PCK@0.05} & \textbf{PCK@0.1} & \textbf{PCK@0.2} \\
    \midrule
    Layer 9  & 9.22\%  & 28.66\% & 53.12\% \\
    Layer 10 & 10.18\% & 31.41\% & 57.77\% \\
    Final    & \textbf{10.60\%} & \textbf{32.68\%} & \textbf{60.13\%} \\
    \bottomrule
    \end{tabular}
\end{table}

%-------------------------------------------------------------------------

\subsubsection{Quantitative Results}

Table~\ref{tab:main_results} reports the quantitative results of fine-tuned models on SPair-71k.
DINO-based models substantially outperform SAM across all PCK thresholds.
DINOv3 achieves the highest overall performance, with $75.77\%$ PCK@0.1 and $86.23\%$ PCK@0.2, while DINOv2 performs slightly better at the strictest threshold (PCK@0.05).

\begin{table}
    \centering
    \setlength{\tabcolsep}{2pt}
    \begin{tabular}{@{}lcccc@{}}
        \toprule
        Model & PCK@0.05 & PCK@0.1 & PCK@0.15 & PCK@0.2 \\
        \midrule
        SAM Finetuned & 19.23 & 31.87 & 40.81 & 47.87 \\
        DINOv2 Finetuned & \textbf{60.02} & 74.89 & 81.39 & 85.37 \\
        DINOv3 Finetuned & 58.15 & \textbf{75.77} & \textbf{82.34} & \textbf{86.23} \\
        \bottomrule
    \end{tabular}
    \caption{\textbf{Per-keypoint PCK@0.1 results on SPair-71k for fine-tuned models.}}
    \label{tab:main_results}
\end{table}

%-------------------------------------------------------------------------

\paragraph{Effect of Fine-Tuning}

Fine-tuning leads to consistent improvements across all models.
For SAM, PCK@0.1 increases from $23.24\%$ to $31.87\%$, demonstrating that task-specific supervision is beneficial.
However, even after fine-tuning, SAM remains significantly behind DINO-based models, indicating architectural limitations for dense correspondence.

%-------------------------------------------------------------------------

\subsubsection{Category-Level Analysis}

We report per-category correspondence accuracy in Table~\ref{tab:category_pck} using PCK@0.1 to analyze model behavior across different object types.
Overall, DINO-based models achieve consistently strong performance across nearly all categories, substantially outperforming SAM by a large margin.
The gap is particularly evident for articulated and deformable objects such as \emph{bird}, \emph{cat}, \emph{dog}, \emph{horse}, and \emph{person}, where DINO models frequently exceed 75--90\% PCK@0.1, while SAM remains below 50\%.

Rigid object categories such as \emph{train}, \emph{bus}, and \emph{car} also benefit from DINO representations, indicating that the learned features capture both semantic and structural regularities.
While DINOv2 and DINOv3 exhibit comparable mean performance, category-specific differences can be observed.
DINOv2 tends to perform better on several articulated categories, whereas DINOv3 achieves higher accuracy on visually distinctive or texture-rich classes such as \emph{train}, \emph{tvmonitor}, and \emph{pottedplant}.
These complementary strengths are consistent with the global trends observed across PCK thresholds.

In contrast, SAM shows limited effectiveness for fine-grained semantic correspondence, with particularly low performance on small objects and categories with high intra-class variability.
This suggests that, despite its strong general-purpose representations, SAM features are less directly suited for precise keypoint localization without more extensive adaptation.

\begin{table}[t]
\centering
\small
\setlength{\tabcolsep}{5pt}
\begin{tabular}{lccc}
\toprule
\textbf{Category} & \textbf{DINOv2} & \textbf{DINOv3} & \textbf{SAM} \\
\midrule
aeroplane     & 81.70 & 74.23 & 31.26 \\
bicycle       & 65.96 & 65.91 & 22.92 \\
bird          & 90.73 & 88.55 & 50.48 \\
boat          & 55.88 & 53.87 & 21.15 \\
bottle        & 60.88 & 64.53 & 32.20 \\
bus           & 85.04 & 77.31 & 27.47 \\
car           & 75.36 & 73.47 & 27.78 \\
cat           & 88.14 & 84.62 & 45.00 \\
chair         & 56.68 & 50.55 & 22.81 \\
cow           & 86.47 & 86.13 & 34.83 \\
dog           & 80.67 & 76.66 & 27.32 \\
horse         & 77.38 & 76.45 & 26.98 \\
motorbike     & 72.81 & 69.93 & 22.63 \\
person        & 78.26 & 70.83 & 31.17 \\
pottedplant   & 48.89 & 68.80 & 27.77 \\
sheep         & 73.47 & 73.03 & 21.21 \\
train         & 88.57 & 90.36 & 40.85 \\
tvmonitor     & 63.31 & 85.31 & 34.81 \\
\midrule
\textbf{Mean} & \textbf{71.48} & \textbf{71.91} & \textbf{27.71} \\
\bottomrule
\end{tabular}
\caption{
\textbf{Per-category PCK@0.1 results on SPair-71k for fine-tuned models.}
DINO-based models substantially outperform SAM across all categories, while DINOv2 and DINOv3 exhibit complementary strengths depending on object type.
}
\label{tab:category_pck}
\end{table}

%-------------------------------------------------------------------------

\subsection{Discussion}

Our results highlight that self-supervised vision transformers trained for representation learning are inherently well-suited for semantic correspondence.
DINOv2 provides strong localization accuracy, while DINOv3 demonstrates improved robustness at higher thresholds, suggesting better global alignment.
Although fine-tuning improves SAM, its performance gap remains substantial, emphasizing that segmentation-oriented architectures are less effective for correspondence tasks.





%-------------------------------------------------------------------------

\subsection{Rerun evaluations with soft-argmax enabled}

%-------------------------------------------------------------------------

\subsection{Generalization across dataset PF-Willow}

